\documentclass[a4paper]{scrartcl}
\usepackage[
    typ=lzk,
    klausurtyp=kurs,
    klasse=scrartcl,
    fach=LF05 Physik,
    lerngruppe=T19E,
    loesungen=seite,
    module={Symbole,Bewertung}    
    %erwartungshorizontAnzeigen,
    %erwartungshorizontStil=simpel,
hHh]{schule}
% Dieses Dokument gehört zu den Beispiel des LaTeX Paketes Schule und ist von den Autoren
% des Pakets erstellt worden.
%
% Das Dokument steht unter der Lizenz: Creative Commons by-nc-sa Version 4.0
% http://creativecommons.org/licenses/by-nc-sa/4.0/deed.de
%
% Nach dieser Lizenz darf das Dokument beliebig kopiert und bearbeitet werden,
% sofern das Folgeprodukt wiederum unter gleichen Lizenzbedingungen vertrieben
% und auf die ursprünglichen Urheber verwiesen wird.
% Eine kommerzielle Nutzung ist ausdrücklich ausgeschlossen.

\author{Heiko Schröter}
\date{\today}
\title{Dynamik\\
(Arbeit, Leistung, Impuls und Stoß)}
\subtitle{}

\usepackage{amsmath}
\usepackage{pgfplots}
\usepackage{mathrsfs}
\usepackage{xfrac}
\usepackage{siunitx}
%\sisetup{quotient-mode = fraction, locale = DE}
\sisetup{per-mode = fraction, locale = DE}
\usepackage{blindtext}
\usepackage{enumitem}
\newcommand{\Ergebnis}[1]{\underline{\underline{#1}}}

\begin{document}
\punktuebersicht*

\setzeAufgabentemplate{schule-randpunkte} 

\begin{aufgabe}[points={4}]
	Die in einem Ventil eingebaute Rückstellfeder mit der Federkonstante $D=\SI{20000}{\newton\per\meter}$ ist mit der Energie $E_{pot}=\SI{25}{\joule}$ gespannt.\\
	\begin{enumerate}[label=(\alph*)]
  	\item Wie groß ist der Federweg $s$, um den die Feder zusammengedrückt ist?
  	\item Wie groß ist die dazu aufzuwendende Spannkraft?
	\end{enumerate}
	\begin{tikzpicture}
		\draw[step=0.5cm,gray,very thin] (0,0) grid (15,8);
	\end{tikzpicture}
	
    \begin{loesung}
    \begin{enumerate}[label=(\alph*)]
  \item
  \begin{flalign*}
  E_{pot}&=W_S=\dfrac{1}{2}\cdot D\cdot s^{2}\Rightarrow s=\sqrt{\dfrac{2\cdot E_{pot}}{D}}\\
  s&=\sqrt{\dfrac{2\cdot \SI{25}{\newton\meter}}{\SI{20000}{\newton\per\meter}}}\\
  &=\SI{0,05}{\meter}=\Ergebnis{\SI{5}{\centi\meter}}&\\
  \end{flalign*}
  \item
    \begin{flalign*}
	F&=D\cdot s=\SI{20000}{\newton\per\meter}\cdot\SI{0,05}{\meter}\\
	&=\Ergebnis{\SI{1000}{\newton}}&\\
  \end{flalign*}
	\end{enumerate}
    \end{loesung}
\end{aufgabe}
\vspace{0.3cm}
\newpage
\begin{aufgabe}[points={5}]
	Eine Kugel fällt aus einer maximalen Höhe von \SI{17,2}{\centi\meter}.\\
	\begin{enumerate}[label=(\alph*)]
  	\item Mit welcher Geschwindigkeit schlägt die Kugel beim ersten Aufschlag auf dem Boden auf?
  	\item  Welche Geschwindigkeit hat sie nach dem ersten Aufschlag \SI{10}{\centi\meter} über dem Boden?
	\end{enumerate}
	\begin{tikzpicture}
		\draw[step=0.5cm,gray,very thin] (0,0) grid (15,8);
	\end{tikzpicture}
	
    \begin{loesung}
    \begin{enumerate}[label=(\alph*)]
  \item
  \begin{flalign*}
	E_{pot}&=E_{kin}\quad\Rightarrow\quad m\cdot g\cdot h=\frac{1}{2}mv^2\\
	v&=\sqrt{2\cdot g\cdot h}=\sqrt{2\cdot \SI{9,81}{\meter\per\square\second}\cdot \SI{0,172}{\meter}}\\
	&=\SI{1,84}{\meter\per\second}=\Ergebnis{\SI{6,61}{\kilo\meter\per\hour}}&\\
  \end{flalign*}
  \item
    \begin{flalign*}
	v_x&=\sqrt{2\cdot g\cdot\Delta h}=\sqrt{2\cdot \SI{9,81}{\meter\per\square\second}\cdot\SI{0,072}{\meter}}\\
	&=\SI{1,19}{\meter\per\second}=\Ergebnis{\SI{4,28}{\kilo\meter\per\hour}}&
  \end{flalign*}
	\end{enumerate}
    \end{loesung}
\end{aufgabe}
\vspace{0.3cm}

\begin{aufgabe}[points={2}]
	Ein Aufzug soll mit der Kraft von \SI{1500}{\newton} angetrieben werden und dabei eine Geschwindigkeit von \SI{1,5}{\meter\per\second} haben. Berechnen Sie die erforderliche Leistung.\\
	\begin{tikzpicture}
		\draw[step=0.5cm,gray,very thin] (0,0) grid (15,4);
	\end{tikzpicture}
	
    \begin{loesung}
	\begin{flalign*}
	&P=F_s\cdot v=\SI{1500}{\newton}\cdot\SI{1,5}{\meter\per\second}=\SI{2250}{\newton\meter\per\second}=\SI{2,25}{\kilo\watt}&
	\end{flalign*}
    \end{loesung}
\end{aufgabe}
\vspace{0.3cm}
\newpage
\begin{aufgabe}[points={4}]
	Ein PKW mit der Masse \SI{1300}{\kilogram} und der Geschwindigkeit \SI{54}{\kilo\meter\per\hour} erhält von einem von hinten auffahrenden Wagen einen Kraftstoß von \SI{2600}{\newton\second}.
	\begin{enumerate}[label=(\alph*)]
	\item Wie groß ist seine Geschwindigkeit unmittelbar nach dem Unfall?
	\item Welchen Kraftstoß erfuhr ein Insasse mit der Masse \SI{68}{\kilogram}?
	\item Welche Kraft wirkte auf das Fahrzeug, wenn es \SI{0,5}{\second} lang beschleunigt worden ist?
	\end{enumerate}
	\begin{tikzpicture}
		\draw[step=0.5cm,gray,very thin] (0,0) grid (15,8);
	\end{tikzpicture}
	
    \begin{loesung}
    \begin{flalign*}
    &v_v=\SI{54}{\kilo\meter\per\hour}=\SI{15}{\meter\per\second}&
    \end{flalign*}
    \begin{enumerate}[label=(\alph*)]
	\item
		\begin{flalign*}
		&F\cdot \Delta t=m\cdot v_n-m\cdot v_v=m\cdot(v_n-v_v)\Rightarrow\\
		&v_n-v_v=F\cdot\dfrac{\Delta t}{m}=\dfrac{\SI{2600}{\newton\second}}{\SI{1300}{\kilogram}}=\SI{2,0}{\meter\per\second}\Rightarrow\\
		&v_n=v_v+\SI{2,0}{\meter\per\second}=\SI{15}{\meter\per\second}+\SI{2,0}{\meter\per\second}=\SI{17}{\meter\per\second}=\Ergebnis{\SI{61,2}{\kilo\meter\per\hour}}&
		\end{flalign*}
	\item    
		\begin{flalign*}
		&F\cdot\Delta t=m\cdot(v_2-v_v)=\SI{68}{\kilogram}\cdot\SI{2}{\meter\per\second}=\Ergebnis{\SI{136}{\newton\second}}&
		\end{flalign*}
	\item
		\begin{flalign*}
		&F=a\cdot m=\dfrac{v_n-v_v}{\Delta t}\cdot m=\dfrac{\SI{2}{\meter\per\second}}{\SI{0,5}{\second}}\cdot \SI{1300}{\kilogram}=\Ergebnis{\SI{5200}{\newton}}\\
		&\text{oder}\quad F=\dfrac{F\cdot\Delta t}{\Delta t}=\dfrac{\SI{2600}{\newton\second}}{\SI{0,5}{\second}}=\Ergebnis{\SI{5200}{\newton}}&
		\end{flalign*}	
	\end{enumerate}
    \end{loesung}
\end{aufgabe}
\vspace{0.3cm}
\newpage
\begin{aufgabe}[points={3}]
Beim Schmieden eines Werkstückes wird ein \SI{1500}{\gram} Hammer verwendet. Die Hammergeschwindigkeit beim Auftreffen beträgt \SI{6}{\meter\per\second}, und der Schmiedeweg (Maß der Werkstückdeformation) beträgt \SI{1.5}{\milli\meter}. Dies bedeutet, dass der Hammer auf einem Weg von \SI{1,5}{\milli\meter} vollkommen abgebremst wird. Berechnen Sie die mittlere Hammerkraft.\\
	\begin{tikzpicture}
		\draw[step=0.5cm,gray,very thin] (0,0) grid (15,6);
	\end{tikzpicture}
	
    \begin{loesung}    
		\begin{flalign*}
		&F\cdot t=m\cdot \Delta v\\
		&\text{gleichmäßig beschleunigte Bewegung: }s=\dfrac{\Delta v\cdot t}{2}\rightarrow t=\dfrac{2s}{\Delta v}\\
		\rightarrow &F\cdot\dfrac{2s}{\Delta v}=m\cdot \Delta v\rightarrow F=\dfrac{m\cdot\Delta v^{2}}{2\cdot s}&\\
		&F=\dfrac{\SI{1,5}{\kilogram}\cdot(\SI{6}{\meter\per\second})^{2}}{2\cdot\SI{0.0015}{\meter}}=\Ergebnis{\SI{18000}{\newton}}&
		\end{flalign*}
    \end{loesung}
\end{aufgabe}
\vspace{0.3cm}

\end{document}